\section{Related Work}
\label{sec:related}

\paragraph{Vietnamese VQA and AutoViVQA.}
Vietnamese VQA has moved from extractive datasets (ViVQA, OpenViVQA, ViTextVQA)
toward free-form generative answering. \textbf{AutoViVQA}~\cite{autovivqa}
provides 19{,}411 images and 37{,}077 questions, each with five free-form answers
and a reasoning-type label, and no public test split; published results span
seq2seq baselines, a fine-tuned \textbf{Vintern-1B}, proprietary LLMs, and
\textbf{ViMoE-VQA}.

\paragraph{Frozen-backbone VLMs and the vision--language bridge.}
A large body of work adapts a frozen image encoder and frozen language model by
training only a small connector --- \textbf{Frozen}~\cite{frozen},
\textbf{BLIP-2}~\cite{blip2} (a Q-Former), \textbf{LLaVA}~\cite{llava} (an MLP
projector, LLM un-frozen). Vintern-1B follows the LLaVA recipe but, downstream,
fully fine-tunes the encoder and projector and LoRA-tunes the decoder; we keep
everything frozen except a lightweight bridge and study a five-design capacity
ladder. ``Inference-Optimal VLMs''~\cite{inferenceoptimalvlms} argues a smaller
LLM with more visual tokens often beats a larger LLM with fewer; our ablations
probe the frozen-backbone version and find the frozen decoder, not the visual
token budget, to be the binding constraint.

\paragraph{Parameter-efficient adaptation, and where in the decoder.}
Parameter-efficient methods insert a few trainable weights ---
\textbf{adapters}~\cite{adapters}, \textbf{prefix / prompt
tuning}~\cite{prefixtuning}, \textbf{LoRA}~\cite{lora}. Two questions the LoRA
literature leaves partly open, and that we address for a frozen-backbone VLM, are
\emph{where} in the decoder the adapter goes and how its benefit interacts with
the connector. Interpretability work casts feed-forward layers as key--value
memories~\cite{ffnkvmemory} and attention as inter-position routing, and
adapter-placement studies find attention-only adaptation often competitive with
adapting all linear layers; our result --- an attention LoRA lifts token-F1 while
a same-rank feed-forward LoRA destabilises training --- is a concrete data point
for frozen 0.5\,B decoders on low-resource generative VQA.

\paragraph{Mixture-of-experts and ``reasoning-aware'' routing.}
\textbf{ViMoE-VQA}~\cite{vimoevqa} improves AutoViVQA accuracy by building a new
generative model (frozen CLIP and PhoBERT, a four-expert Top-2 MoE, a six-layer
decoder, trained from scratch) and attributes part of its gain to
``approximate reasoning-aware expert selection''. We are a direct counterpoint on
the \emph{same benchmark}: we improve the existing Vintern-1B rather than replace
it, and we test the reasoning-aware premise directly
(Sec.~\ref{sec:abl-routing}) --- reasoning type carries no signal about the
optimal visual-compute action, and no learned policy beats a fixed one. ViMoE's
own leave-one-out ablation is consistent (removing any expert moves BLEU by only
$0.11$--$0.44$), suggesting the gain is added capacity rather than specialisation.
